%% =====================================================================
%%  T-21-03  Strategic Disdain
%%  -------------------------------------------------------------------
%%  One idea per file. Core is mandatory. Slots in canonical order.
%%  Max 700 words. No layout commands. No filler. New source material
%%  for this entry goes in sources/<BOOK>/T-21-03.tex -- never by
%%  rewriting this file (docs/RULES.md R0.1).
%%  =====================================================================
%% @kind: topic
%% @id: T-21-03
%% @title: Strategic Disdain
%% @subtitle: Pricing the unattainable as beneath you
%% @chapter: c21
%% @order: 30
%% @status: stable
%% @sources: B3
%% @updated: 2026-09-19
%% @allow-rewrite: B3 ingest 2026-09-19 -- committed scaffold replaced by the written entry (planned -> stable lifecycle)

\topic{T-21-03}{Strategic Disdain}{Pricing the unattainable as beneath you}

\begin{core}
Unattainable things are prizes; that is the market's default pricing. Disdain is the counter-move: publicly refusing to want what you cannot have converts the unattainable into the irrelevant, and --- the deeper mechanism --- a deliberate grain of disdain sprinkled into your dealings makes your regard itself scarce and valuable. To disregard is to win regard.
\end{core}
\begin{mechanism}
Desire is visible, and visible desire is leverage: whoever can tell what you want can price it out of your reach or sell it back to you at passion's premium. Public disdain removes the lever --- the prize you will not chase cannot be dangled. The second mechanism runs on scarcity of approval: most people distribute warmth freely, which makes their warmth worthless; the person whose regard must be earned, who can conspicuously dispense with your company, raises the price of their own attention. The source's advice is almost comical in its precision: let those who depend on you feel, from time to time, that they could do without your company --- a grain of disdain mixed into your dealings, Schopenhauer says, makes your friendship worth more, and even a dog cannot endure being treated too kindly. The mechanism underneath is standard scarcity: humans value what others demonstrably refuse.
\end{mechanism}
\begin{conditions}
Strongest where attention is the currency --- courts, scenes, institutions with more suitors than seats --- and in any negotiation where your wanting has leaked. Weakens among secure peers who read indifference as information rather than status, in genuine partnerships where disdain is simply damage, and on anyone immune to your opinion, where the grain of disdain is just rudeness with a strategy's name attached.
\end{conditions}
\begin{application}
  \item Inventory your visible wants. Anything you cannot afford to have, stop being seen to want --- remove the tells, not just the claim.
  \item Exit one chase publicly, gracefully, without commentary. The unexplained renunciation does the work.
  \item With dependents and suitors, ration availability: warm after work, cool at request. The pattern is the message.
  \item Let one invitation lapse unchased per season. Scarcity needs evidence, not announcements.
  \item Never disdain what you visibly still want --- the gap between claim and behaviour is where credibility dies.
\end{application}
\begin{feedback}
  \item Working: the unattainable thing's price --- or its gatekeepers --- starts courting you. Refusal re-priced it.
  \item Working: your attention gets treated as an event; invitations arrive earlier and better.
  \item Failing: the room reads petulance --- disdain aimed at something everyone knows you wanted.
  \item Failing: the dependents stop bringing you information. You priced your attention past its utility.
\end{feedback}
\begin{failure}
Disdain is a performance with a morale cost: performed long enough, it teaches the performer to want nothing, which is a kind of impoverishment no balance sheet shows. It also selects badly --- used on the insecure it manufactures resentment rather than regard, and resentment waits. The strategic version also collapses on contact with genuine need: disdain for what you truly cannot survive without is not pricing, it is self-deception with good posture.
\end{failure}
\begin{countermeasures}
  \item Notice who is rationing their regard around you. Ask what they get paid in --- scarcity performances have a buyer in mind, and it is usually you.
  \item Price attention by content, not difficulty. Hard-to-get warmth is still warmth you did not need.
  \item Audit your own chases: what are you pursuing mainly because it is withheld? That is the reactance trigger wearing your face (T-10-02).
  \item Give real regard freely to the people you keep. Its value to them will outprice any rival's theatre.
\end{countermeasures}

\sources{B3}
\seesources
\seealso{T-10-02, T-21-01, T-22-01}
